\documentclass[10pt,a4paper]{amsart}
\usepackage{amsmath,amssymb,amsthm,graphicx}
% Retained preamble for the Stacks Project chapter morphisms.tex (Morphisms of Schemes).
% The chapter file pinned-original-source.tex carries only \input{preamble} on its first line, so no
% preamble text lives inside the chapter itself.  The Stacks Project uses one single shared preamble for
% every chapter; the snapshot stages it verbatim in sources/stacks-chow/preamble.tex.  The block below is
% copied mechanically, line for line, from that file's lines 14-18 (the Xy-pic setup) and lines 154-217
% (the native theorem declarations together with the mathematical macros).  Nothing was added, removed or
% redefined.  The wrapper class and the amsmath/amssymb/amsthm/graphicx base packages are supplied by the
% assembler, so the class/\externaldocument part of the shared preamble is not repeated here.
% For commutative diagrams we use Xy-pic
\usepackage[all]{xy}

% We use 2cell for 2-commutative diagrams.
\xyoption{2cell}
% Theorem environments.
%
\theoremstyle{plain}
\newtheorem{theorem}[subsection]{Theorem}
\newtheorem{proposition}[subsection]{Proposition}
\newtheorem{lemma}[subsection]{Lemma}

\theoremstyle{definition}
\newtheorem{definition}[subsection]{Definition}
\newtheorem{example}[subsection]{Example}
\newtheorem{exercise}[subsection]{Exercise}
\newtheorem{situation}[subsection]{Situation}

\theoremstyle{remark}
\newtheorem{remark}[subsection]{Remark}
\newtheorem{remarks}[subsection]{Remarks}

\numberwithin{equation}{subsection}

% Macros
%
\def\lim{\mathop{\mathrm{lim}}\nolimits}
\def\colim{\mathop{\mathrm{colim}}\nolimits}
\def\Spec{\mathop{\mathrm{Spec}}}
\def\Hom{\mathop{\mathrm{Hom}}\nolimits}
\def\Ext{\mathop{\mathrm{Ext}}\nolimits}
\def\SheafHom{\mathop{\mathcal{H}\!\mathit{om}}\nolimits}
\def\SheafExt{\mathop{\mathcal{E}\!\mathit{xt}}\nolimits}
\def\Sch{\mathit{Sch}}
\def\Mor{\mathop{\mathrm{Mor}}\nolimits}
\def\Ob{\mathop{\mathrm{Ob}}\nolimits}
\def\Sh{\mathop{\mathit{Sh}}\nolimits}
\def\NL{\mathop{N\!L}\nolimits}
\def\CH{\mathop{\mathrm{CH}}\nolimits}
\def\proetale{{pro\text{-}\acute{e}tale}}
\def\etale{{\acute{e}tale}}
\def\QCoh{\mathit{QCoh}}
\def\Ker{\mathop{\mathrm{Ker}}}
\def\Im{\mathop{\mathrm{Im}}}
\def\Coker{\mathop{\mathrm{Coker}}}
\def\Coim{\mathop{\mathrm{Coim}}}

% Boxtimes
%
\DeclareMathSymbol{\boxtimes}{\mathbin}{AMSa}{"02}

%
% Macros for moduli stacks/spaces
%
\def\QCohstack{\mathcal{QC}\!\mathit{oh}}
\def\Cohstack{\mathcal{C}\!\mathit{oh}}
\def\Spacesstack{\mathcal{S}\!\mathit{paces}}
\def\Quotfunctor{\mathrm{Quot}}
\def\Hilbfunctor{\mathrm{Hilb}}
\def\Curvesstack{\mathcal{C}\!\mathit{urves}}
\def\Polarizedstack{\mathcal{P}\!\mathit{olarized}}
\def\Complexesstack{\mathcal{C}\!\mathit{omplexes}}
% \Pic is the operator that assigns to X its picard group, usage \Pic(X)
% \Picardstack_{X/B} denotes the Picard stack of X over B
% \Picardfunctor_{X/B} denotes the Picard functor of X over B
\def\Pic{\mathop{\mathrm{Pic}}\nolimits}
\def\Picardstack{\mathcal{P}\!\mathit{ic}}
\def\Picardfunctor{\mathrm{Pic}}
\def\Deformationcategory{\mathcal{D}\!\mathit{ef}}

% --- appended 2026-10-12 (Chief): the project own definitions of the reference/slogan/history
% display environments, copied mechanically from the shared Stacks preamble lines 8-12
% (\usepackage{verbatim} plus the three \newenvironment lines, which the chapter preamble keeps
% verbatim from the project preamble). Nothing added, removed or redefined: these environments
% carry no mathematical content and the project itself typesets them as comments. Without these
% three lines a counted statement whose own text uses them fails with Environment slogan undefined.
\usepackage{verbatim}
\newenvironment{reference}{\comment}{\endcomment}
\newenvironment{slogan}{\comment}{\endcomment}
\newenvironment{history}{\comment}{\endcomment}
\usepackage{cleveref}
\crefname{originallemma}{original lemma}{original lemmas}
\Crefname{originallemma}{Original lemma}{Original lemmas}
\makeatletter
\newlabel{morphisms-lemma-seminormal-local-property}{{0EUM}{source}{}{sourcepointer-1}{}}
\newlabel{morphisms-lemma-seminormal-local-property@cref}{{[originallemma][0][]0EUM}{[1][0][]source}{}{}{}}
\makeatother
\crefname{originallemma}{original lemma}{original lemmas}
\Crefname{originallemma}{Original lemma}{Original lemmas}
\makeatletter
\newlabel{properties-definition-property-local}{{01OP}{source}{}{sourcepointer-2}{}}
\newlabel{properties-definition-property-local@cref}{{[originallemma][0][]01OP}{[1][0][]source}{}{}{}}
\makeatother
\crefname{originallemma}{original lemma}{original lemmas}
\Crefname{originallemma}{Original lemma}{Original lemmas}
\makeatletter
\newlabel{algebra-lemma-standard-covering}{{00EJ}{source}{}{sourcepointer-3}{}}
\newlabel{algebra-lemma-standard-covering@cref}{{[originallemma][0][]00EJ}{[1][0][]source}{}{}{}}
\makeatother
\title{Seminormality and absolute weak normality are local properties of rings}
\author{The Stacks Project Authors}
\date{stacks-project:morphisms.tex:local-snapshot-20261012; GNU FDL 1.2 or later}
\begin{document}
\maketitle
\setcounter{subsection}{1}
\noindent\textit{Editorial note.} All delivered mathematical text is the author's own, copied from the pinned Stacks Project chapter file morphisms.tex; no statement or proof step was rewritten, reordered or removed.\par\medskip
\noindent\textit{Editorial note.} The chapter file pinned-original-source.tex carries only \textbackslash{}input\{preamble\} on its first line, so no preamble text lives inside the chapter. The retained file sources/stacks-morphisms/retained-preamble.tex is a mechanical, line-for-line copy of the project's single shared preamble as staged in the snapshot. Nothing was added, removed or redefined, and the excerpt serves only as preamble, never as a span coordinate.\par\medskip
\noindent\textit{Editorial note.} Both counted spans start at the \textbackslash{}begin line of their environment, so the statement and the author's proof are delivered whole; surrounding prose of the section is deliberately outside the counted spans and is carried as standard prerequisites instead.\par\medskip
\noindent\textit{Editorial note.} The citation fields (source\_id, source\_version, authors, source\_url, license.name, source\_title) are identical to those of every other card this lane has produced from this chapter, so the batch selector treats them as one source.\par\medskip
\noindent\textit{Editorial note.} Screened under the Lead's 12:4x rule (prescreen hit, proof body 1136 bytes): read in full and accepted, not refused. The seminar-normal half is proved out in detail; the absolutely weakly normal half is dispatched by the author's own words 'In exactly the same manner the reader shows that if \$A\$ is absolutely weakly normal, then \$A\_f\$ is absolutely weakly normal' and 'then the exact same argument shows that \$A\$ is absolutely weakly normal' -- a self-description of a routine repetition of the two arguments given just above, which is the class the Lead's rule says to accept with the wording quoted. It is deliberately NOT in the class of the three refusals in this same window (lemma-seminormalization-ring, lemma-seminormalization, lemma-rational-map-from-reduced-to-separated), where a numbered statement item is left with no argument whatsoever ('we omit the proof of (3) and (4)'). Both verbatim phrases are quoted here so that a reviewer can re-judge the call without re-reading the source.\par\medskip
\noindent\textit{Allowed standard prerequisites (editorial).}\begin{itemize}
\item The definition of a local property of rings, in terms of localisations at a finite family generating the unit ideal. Properties, Definition \textbackslash{}ref\{properties-definition-property-local\}, cited at proof lines 12023-12026. Fixes what has to be proved for the converse direction.
\item Glueing of elements given over a standard covering of the spectrum. Algebra, Lemma \textbackslash{}ref\{algebra-lemma-standard-covering\}, cited at proof lines 12047-12050. Produces the single element inducing all the local ones.
\item The definition of seminormal and absolutely weakly normal rings. Definition \textbackslash{}ref\{definition-seminormal-ring\}, the definition that this lemma is about. Supplies the equations and the uniqueness clause used in both directions.
\end{itemize}\medskip
\ifdefined\hypertarget\hypertarget{sourcepointer-1}{}\fi
\noindent\textit{Source reference.} Original lemma 0EUM: focused-seeds/AGP/sources/stacks-snapshot-20261011/tags.txt, source lines 18557--18557. The lemma's own canonical tag, quoted so that readers can resolve the label inside the delivered text; it points at the counted statement itself and carries no proof content.\par\medskip
\ifdefined\hypertarget\hypertarget{sourcepointer-2}{}\fi
\noindent\textit{Source reference.} Original lemma 01OP: focused-seeds/AGP/sources/stacks-snapshot-20261011/tags.txt, source lines 2202--2202. Navigation pointer to a statement the counted proof refers to but does not contain; nothing load-bearing is delegated to it.\par\medskip
\ifdefined\hypertarget\hypertarget{sourcepointer-3}{}\fi
\noindent\textit{Source reference.} Original lemma 00EJ: focused-seeds/AGP/sources/stacks-snapshot-20261011/tags.txt, source lines 536--536. Navigation pointer to a statement the counted proof refers to but does not contain; nothing load-bearing is delegated to it.\par\medskip
\begin{lemma}
\label{lemma-seminormal-local-property}
Being seminormal or being absolutely weakly normal
is a local property of rings, see
Properties, Definition \ref{properties-definition-property-local}.
\end{lemma}

\begin{proof}
Suppose that $A$ is seminormal and $f \in A$. Let $x', y' \in A_f$
with $(x')^3 = (y')^2$. Write $x' = x/f^{2n}$ and $y' = y/f^{3n}$
for some $n \geq 0$ and $x, y \in A$. After replacing $x, y$
by $f^{2m}x, f^{3m}y$ and $n$ by $n + m$, we see that
$x^3 = y^2$ in $A$. Then we find a unique $a \in A$ with
$x = a^2$ and $y = a^3$. Setting $a' = a/f^n$ we get
$x' = (a')^2$ and $y' = (a')^3$ as desired. Uniqueness
of $a'$ follows from uniqueness of $a$.
In exactly the same manner the reader shows that if $A$ is
absolutely weakly normal,
then $A_f$ is absolutely weakly normal.

\medskip\noindent
Assume $A$ is a ring and $f_1, \ldots, f_n \in A$ generate the unit
ideal. Assume $A_{f_i}$ is seminormal for each $i$.
Let $x, y \in A$ with $x^3 = y^2$. For each $i$ we find a unique
$a_i \in A_{f_i}$ with $x = a_i^2$ and $y = a_i^3$ in $A_{f_i}$.
By the uniqueness and the result of the first paragraph (which
tells us that $A_{f_if_j}$ is seminormal) we see that
$a_i$ and $a_j$ map to the same element of $A_{f_if_j}$.
By Algebra, Lemma \ref{algebra-lemma-standard-covering}
we find a unique $a \in A$
mapping to $a_i$ in $A_{f_i}$ for all $i$.
Then $x = a^2$ and $y = a^3$ by the same token. Clearly this $a$ is unique.
Thus $A$ is seminormal. If we assume $A_{f_i}$ is absolutely weakly normal,
then the exact same argument shows that $A$ is absolutely weakly normal.
\end{proof}

\end{document}
